\chapter*{使用说明}
\addcontentsline{toc}{chapter}{使用说明}

本讲义采用经典数值分析教材的知识组织方式：章节标题使用稳定的学科术语，问题驱动的解释放在章节内部。阅读顺序遵循“现象与算法结构 $\rightarrow$ 数学形式化 $\rightarrow$ 收敛分析 $\rightarrow$ 并行执行 $\rightarrow$ 软件实现与性能诊断”。因此，局部 Fourier 分析（LFA）被安排在读者已经完整走过 two-grid 与 V-cycle 之后，而不是作为第一次接触多重网格时的前置门槛。

全书分为四部分。第一部分先建立单层椭圆型离散，再把它推广到自适应网格细化（adaptive mesh refinement, AMR）的复合离散系统，随后研究平滑迭代、粗网格校正、V-cycle 与 LFA；第二部分先建立架构无关的并行计算模型和多重网格操作级并行模型，再把同一套 task/data/ownership 语义分别映射到共享内存、GPU、MPI、多 GPU 与可移植后端；第三部分进入多重网格软件体系、性能评测与瓶颈诊断；第四部分讨论低精度计算、专用矩阵硬件、matrix-free 与粗网格通信等现代方向。

离散部分沿两条并行路线组织：有限差分（FD）从微分算子的离散差分出发，有限体积（FV）从控制体积分与数值通量出发。二者描述的是不同的构造视角，存在重叠而非互斥；在规则笛卡尔网格的守恒格式中，同一个离散算子可以同时具有 flux-difference FD 与 FV 的解释。\chapref{ch:discretization}负责建立两条路线及其交集，\chapref{ch:amr-discretization}继续区分 cell-centered FV、守恒型 FD 与节点型 FD 在粗细接口上的语义。自\chapref{ch:iterative-smoothing}起，主线提升到代数算子 $A$ 与 multigrid hierarchy；只有当结论依赖局部守恒、对称性、stencil 几何、系数位置或变量语义时，正文才回到具体的 FD/FV 离散来源。

除少数纯推导附录外，各章统一采用“本章问题 $\rightarrow$ 最小例子 $\rightarrow$ 术语定义 $\rightarrow$ 一般公式/算法 $\rightarrow$ 工程含义 $\rightarrow$ 习题”的展开顺序。目录中的章、节、小节只使用稳定的学科术语承担知识分类；问题式表述只出现在正文导语和例题中。高级术语第一次出现时要么立即定义，要么明确标注为后文内容；不使用尚未建立的概念去解释另一个新概念。允许为了真实工程语境提前出现后文术语，但第一次出现必须同时给出\textbf{最低可用定义}和\textbf{明确的后文交叉引用}，让读者知道“现在先理解到什么程度”以及“正式展开在哪里”。此外，所有可能因插章、重排或增删公式而变化的章节、节、公式、图、表和算法编号，一律通过 \texttt{label/ref/eqref} 交叉引用生成，不在正文中手写编号。

跨章概念按“\textbf{预告 $\rightarrow$ 最小模型 $\rightarrow$ 范围边界 $\rightarrow$ 后文定位 $\rightarrow$ 正式回扣}”组织。术语可以提前出现，但第一次出现只给当前推理所需的最低含义，并用交叉引用指出正式展开的位置；到达正式章节后，再把定义、推导和适用范围补全，并明确它与先前语境的关系。这样既保留真实工程语境，也让读者始终知道一个新名词当前需要理解到什么程度、以后在哪里继续学习。

章节内部还遵循三条结构规则。第一，\textbf{维数提升必须重新建模}：一维用于建立最小直觉；进入二维时必须显式说明新增的拓扑、矩阵块结构、边界角点与切向方向；进入三维时还要说明切片、棱、角与两个切向方向，不能用“同理”代替新增结构。第二，\textbf{先定义评价对象，再介绍改变它的方法}：例如进入迭代法之前先区分真解、误差与可计算残差，再讨论 Jacobi 或 Gauss--Seidel 怎样构造修正。第三，\textbf{section 对应稳定的认知阶段，而不是一个公式、一个后端名称或一个很短的变体}；同一认知阶段内的方法差异使用 subsection/paragraph 组织。这样目录本身应当能够读成一条因果链，而不是术语清单。


本书的章末习题固定包含四类训练，但各类比例随章节主题变化：
\begin{enumerate}
\item \textbf{计算题}：采用线性代数、高等数学与经典数值分析的风格，在有限规模对象上手算矩阵、向量、迭代步骤、特征值、离散系数、Galerkin 粗算子、Fourier symbol 等，要求给出中间步骤而不只写结论；
\item \textbf{推导与证明题}：从定义和已建立公式推出一般结论，训练误差分析、正定性、投影性质、复杂度与并行正确性等结构性推理；
\item \textbf{编程与上机题}：把手算的小问题推广到更大规模，通过实现、数值实验、收敛曲线和 profiler 数据验证前两类题中的数学结论；
\item \textbf{工程与综合题}：把数值收敛、数据结构、硬件成本和软件设计放到同一问题中做判断。性能估算属于这一类或并行章节的计算题，但不替代经典数值分析计算训练。
\end{enumerate}
推荐的学习顺序是“\textbf{先手算一个小问题 $\\rightarrow$ 再推导一般式 $\\rightarrow$ 再编程放大验证 $\\rightarrow$ 最后做工程判断}”。四类题都不是可选装饰；它们共同决定读者是否真正获得“会算、会证、会写、会测”的能力。

\begin{keybox}
\textbf{贯穿全文的核心问题：}多重网格在数学上通过尺度分解消除不同频率的误差；并行实现则必须让这些尺度在不同硬件上维持足够的并行度、数据局部性与低同步成本。粗网格阶段是数值算法与并行体系结构最直接发生冲突的位置。
\end{keybox}

\begin{warningbox}
\textbf{术语约定：}\chapref{ch:amr-discretization}中的 level 指 AMR 的物理离散层级；多重网格章节随后会引入求解器自己的 MG level。两类 level 同时出现以后，正文始终显式写成 AMR level 或 MG level，不用裸露的 ``level'' 代替二者。
\end{warningbox}
